\documentclass[12pt]{article}
\usepackage[utf8]{inputenc}
\usepackage[T1]{fontenc}
\usepackage[margin=1in]{geometry}
\usepackage{setspace}
\usepackage{xcolor}
\usepackage[colorlinks=true, urlcolor=blue!60!black, linkcolor=blue!60!black]{hyperref}
\setlength{\parindent}{2em}
\setlength{\parskip}{1.25em}
\renewcommand{\baselinestretch}{1.0}
\begin{document}
\begin{center}
{\LARGE \textit{Trump v. Slaughter}, 609 U.S. 422 (2026)\par}
\vspace{0.8em}
{\large Case Brief\par}
\end{center}
\vspace{1em}

\section{Memory Jogger}

In a landmark 6--3 decision, the Supreme Court overruled \textit{Humphrey's Executor v. United States} (1935) and held that the Federal Trade Commission's statutory ``for-cause'' removal protections violate Article II's separation of powers because principal officers who exercise executive power must be removable by the President at will (\href{https://www.law.cornell.edu/supremecourt/text/25-332}{Cornell LII}).

\section{Detailed Case Facts}

Created by the Federal Trade Commission Act of 1914 (38 Stat. 717, codified at 15 U.S.C. \S~41 \textit{et seq.}), the Federal Trade Commission (FTC) administers and enforces more than 80 federal statutes spanning antitrust, consumer protection, data privacy, and commercial regulation (\href{https://www.supremecourt.gov/opinions/25pdf/25-322_new_4gcj.pdf}{Slip Op. at 2, 25}; \href{https://www.courtlistener.com/opinion/10881681/trump-v-slaughter/}{CourtListener}). The FTC promulgates substantive rules carrying the force of law (including defining ``unfair or deceptive'' acts or practices under 15 U.S.C. \S~57a), investigates businesses and enforces compliance through in-house administrative adjudications, and files civil enforcement suits in federal court seeking injunctions, consumer redress, and billions of dollars in civil penalties. Under 15 U.S.C. \S~41, the FTC is headed by five Commissioners appointed by the President with Senate confirmation for staggered seven-year terms, no more than three of whom may belong to the same political party, and who may be removed by the President only ``for inefficiency, neglect of duty, or malfeasance in office.''

In 2018, during his first term, President Donald J. Trump nominated Democrat Rebecca Kelly Slaughter to serve on the FTC, and the Senate unanimously confirmed her; President Joe Biden renominated her, and she was reconfirmed in 2024 for a second term (\href{https://www.law.cornell.edu/supremecourt/text/25-332}{Cornell LII}; \href{https://www.scotusblog.com/2026/06/court-allows-trump-to-fire-ftc-commissioner-and-overturns-major-restraint-on-presidential-power/}{SCOTUSblog}). When President Trump began his second term in January 2025, the Commission had three Democrats and two Republicans. On his first day in office, Trump designated Republican Andrew Ferguson as Chair, and former Chair Lina Khan soon resigned, leaving an even 2--2 split. In March 2025, President Trump fired the two remaining Democratic Commissioners, Rebecca Slaughter and Alvaro Bedoya. He did not assert inefficiency, neglect of duty, or malfeasance in office; instead, his removal letters stated that their ``continued service on the FTC [was] inconsistent with [his] Administration's priorities'' and invoked his ``authority under Article II of the Constitution'' (\href{https://www.law.cornell.edu/supremecourt/text/25-332}{Cornell LII}). Bedoya later formally resigned from the FTC, leaving Slaughter's claims alone before the courts.

\section{Procedural History}

\begin{itemize}
    \item \textbf{U.S. District Court for the District of Columbia (No. 1:25-cv-00909):} Slaughter (initially joined by Bedoya, whose claims were dismissed as moot after his resignation) sued President Trump and other executive officials, alleging her removal was \textit{ultra vires}, violated the Administrative Procedure Act, and violated the Constitution (\href{https://www.law.cornell.edu/supremecourt/text/25-332}{Cornell LII}). On July 17, 2025, Judge Loren AliKhan granted Slaughter's motion for summary judgment (\textit{Slaughter v. Trump}, 791 F. Supp. 3d 1 (D.D.C. 2025)), holding that \textit{Humphrey's Executor v. United States}, 295 U.S. 602 (1935), remained binding Supreme Court precedent directly on point, declaring the removal \textit{ultra vires}, and issuing a permanent injunction barring interference with Slaughter's duties (\href{https://supreme.justia.com/cases/federal/us/609/25-332/}{Justia}).
    \item \textbf{U.S. Court of Appeals for the D.C. Circuit:} On September 2, 2025, a divided panel (Judges Millett and Pillard) denied the Government's motion for a stay pending appeal (2025 WL 2551247), reasoning that the Government had no prospect of success while \textit{Humphrey's Executor} stood. Judge Rao dissented, contending that \textit{Humphrey's} does not authorize a ``headless fourth branch'' wielding considerable executive power (\href{https://www.law.cornell.edu/supremecourt/text/25-332}{Cornell LII}; \href{https://www.scotusblog.com/2026/06/court-allows-trump-to-fire-ftc-commissioner-and-overturns-major-restraint-on-presidential-power/}{SCOTUSblog}).
    \item \textbf{Supreme Court of the United States (No. 25A264 / No. 25-332):} On September 4, 2025, the Government applied for a stay. On September 22, 2025 (606 U.S. 1051; \href{https://www.supremecourt.gov/opinions/24pdf/25a264_o759.pdf}{Order in 25A264}), over a dissent by Justice Kagan (joined by Justices Sotomayor and Jackson), the Court stayed the District Court's injunction and treated the application as a petition for a writ of certiorari before judgment, directing briefing on (1) whether FTC removal protections violate the separation of powers and whether \textit{Humphrey's Executor} should be overruled, and (2) whether a federal court may prevent a person's removal from public office at equity or law. Oral argument took place on December 8, 2025 (Solicitor General D. John Sauer for petitioners; Amit Agarwal for respondents). On June 29, 2026, the Court reversed and remanded, 6--3 (\href{https://www.supremecourt.gov/opinions/25pdf/25-322_new_4gcj.pdf}{Slip Op. at 36}).
\end{itemize}

\section{Judicial Votes}

\textbf{Vote:} 6--3 reversing and remanding (\href{https://www.law.cornell.edu/supremecourt/text/25-332}{Cornell LII}; \href{https://www.courtlistener.com/opinion/10881681/trump-v-slaughter/}{CourtListener}).

\begin{itemize}
    \item \textbf{Majority (6):} Chief Justice Roberts delivered the opinion of the Court, joined in full by Justices Alito, Gorsuch, Kavanaugh, and Barrett, and joined by Justice Thomas as to all parts except Part III--B.
    \item \textbf{Concurrence (1):} Justice Gorsuch filed a concurring opinion.
    \item \textbf{Dissent (3):} Justice Sotomayor filed a dissenting opinion, joined by Justices Kagan and Jackson.
\end{itemize}

\section{Holding}

\textbf{Judgment:} Reversed and remanded. The District Court's summary judgment and permanent injunction reinstating Commissioner Slaughter are reversed (\href{https://www.supremecourt.gov/opinions/25pdf/25-322_new_4gcj.pdf}{Slip Op. at 36}).

\textbf{Rule Announced:} Article II's Vesting Clause (``[t]he executive Power shall be vested in a President'') and Take Care Clause grant the President the unrestricted authority to remove at will principal officers who exercise executive power on his behalf. Statutory ``for-cause'' removal protections for FTC Commissioners (15 U.S.C. \S~41) violate the constitutional separation of powers, and \textit{Humphrey's Executor v. United States}, 295 U.S. 602 (1935), is overruled to the extent anything remained of it beyond entities exercising ``no part of the executive power'' (\href{https://www.law.cornell.edu/supremecourt/text/25-332}{Cornell LII}).

\textbf{Holding vs. Dicta:} The core holding is that an agency exercising core executive functions---promulgating binding substantive rules, conducting investigations and binding in-house enforcement adjudications, and filing civil lawsuits on behalf of the United States---cannot be led by principal officers insulated from at-will presidential removal, overruling \textit{Humphrey's Executor}. Because Slaughter lost on the constitutional merits (Question 1), the Court did not reach Question 2 (judicial power to order reinstatement), which it addressed instead in \textit{Trump v. Cook}. In Part III--B (which Justice Thomas declined to join), the Court carefully cabined its holding in non-binding reservations/dicta: it declined to define the outer boundaries of ``executive power,'' noted that congressionally chartered entities (such as the Boy Scouts of America or Georgetown University) and purely investigative or informative bodies need not exercise executive power, left open historical exceptions such as the Federal Reserve (decided that same day in \textit{Cook}), and reserved the constitutionality of tenure protections for non-Article III judges (such as the Tax Court and Court of Federal Claims) and questions of statutory severability (\href{https://www.paulweiss.com/insights/client-memos/supreme-court-overrules-humphrey-s-executor-but-preserves-the-fed-s-independence}{Paul, Weiss}).

\section{Analysis of Opinions}

\subsection{Majority Opinion (Chief Justice Roberts, joined by Alito, Gorsuch, Kavanaugh, and Barrett, JJ., and by Thomas, J., as to all but Part III--B)}

Chief Justice Roberts anchored the opinion in constitutional text, structure, and founding-era history. Rejecting the weak, council-checked governorships of post-Revolutionary state constitutions, the Framers deliberately chose a single, unitary President to ensure ``vigor and activity'' and clear democratic accountability---so the people would always know ``on whom the blame or the punishment of a pernicious measure \ldots{} ought really to fall'' (\textit{The Federalist No. 70}; \href{https://www.law.cornell.edu/supremecourt/text/25-332}{Cornell LII}). Because one President cannot execute the laws alone, subordinate officers must serve as his ``assistants or deputies'' subject to his ``superintendence'' (\textit{The Federalist No. 72}), which necessarily requires at-will removal. Roberts traced the liquidation of this principle through the First Congress's ``Decision of 1789,'' the practice of early Presidents from Washington through Jackson, 19th-century decisions (\textit{Ex parte Hennen}, \textit{Parsons v. United States}), and the Court's landmark ruling in \textit{Myers v. United States}, 272 U.S. 52 (1926).

Turning to \textit{Humphrey's Executor}, Roberts argued that the 1935 decision rested on an ``almost fictional view'' that the FTC exercised ``no part of the executive power'' and performed merely ``quasi-legislative'' and ``quasi-judicial'' duties in aid of Congress and the courts. Subsequent decisions---\textit{Buckley v. Valeo}, \textit{INS v. Chadha}, \textit{Bowsher v. Synar}, \textit{Morrison v. Olson}, \textit{Free Enterprise Fund}, and \textit{Seila Law}---progressively dismantled those premises, leaving nothing of \textit{Humphrey's} except the truism that an entity exercising no executive power falls outside \textit{Myers}. Applying the \textit{Knick v. Township of Scott} \textit{stare decisis} factors, the Court formally overruled \textit{Humphrey's}: its reasoning was flawed from inception (``no ifs, ands, or quasis about it''); it conflicted with decades of modern separation-of-powers precedent; its indeterminate ``quasi'' test proved unworkable in lower courts; and congressional reliance on aggrandizing its own power over insulated agencies is not a legitimate reliance interest---``[n]o branch may rely on adverse possession to claim power that the Constitution vests elsewhere'' (\href{https://www.law.cornell.edu/supremecourt/text/25-332}{Cornell LII}).

Applying these principles to the modern FTC, Roberts deemed the case easy: rulemaking, in-house enforcement adjudications, and filing civil suits in federal court are ``executive through and through.'' In Part III--B (which Justice Thomas refused to join, consistent with his dissent in \textit{Cook}), Roberts reserved judgment on non-executive entities, the Federal Reserve, and non-Article III courts. Finally, Roberts rejected Slaughter's proposed ``reasonableness'' test as unprincipled---warning it would allow Congress to insulate Cabinet departments like the DOJ or EPA---and dismissed the dissent's historical examples (the Sinking Fund Commission, Mint Board, and Revolutionary War Debt Commission) as either removable at will or consistent with presidential responsibility (\href{https://supreme.justia.com/cases/federal/us/609/25-332/}{Justia}).

\subsection{Concurring Opinion (Justice Gorsuch)}

Justice Gorsuch joined the majority in full but wrote a striking 16-page solo concurrence warning about the institutional consequences of \textit{Slaughter}. Tracing the Progressive Era administrative theory of Woodrow Wilson and James Landis, Gorsuch observed that Congress delegated sweeping legislative rulemaking, regulatory crime-creation, and private-rights adjudication to independent agencies largely on the assumption that they were insulated from the White House (\href{https://www.law.cornell.edu/supremecourt/text/25-332}{Cornell LII}). Now that \textit{Humphrey's} is gone, ``the fourth branch's powers still exist; they have just been reassigned to the President.'' Warning that a presidential veto makes it nearly impossible for Congress to reclaim delegated authority by simple majority (a ``ratchet effect''), Gorsuch urged the Court to ``finish the journey'' by rigorously enforcing the nondelegation doctrine, the major questions doctrine, vagueness doctrine, Article III, and the Seventh Amendment to restore legislative power to Congress and judicial power to independent courts and juries (\href{https://lawliberty.org/gorsuchs-con-jority-opinion-in-slaughter/}{Law \& Liberty}).

\subsection{Dissenting Opinion (Justice Sotomayor, joined by Kagan and Jackson, JJ.)}

Justice Sotomayor authored a sweeping 49-page dissent arguing that the majority upended 90 years of binding precedent and 140 years of bipartisan political practice---dating back to the 1887 Interstate Commerce Commission and embraced by more than a dozen Presidents---based on an ahistorical theory of illimitable executive power (\href{https://www.law.cornell.edu/supremecourt/text/25-332}{Cornell LII}). Textually, the Constitution is silent on removal outside impeachment, vests office creation in Congress under the Necessary and Proper Clause, and uses the Take Care Clause as a fiduciary duty rather than a grant of boundless power. Historically, neither the English Crown nor state governors held unregulable removal power, \textit{The Federalist No. 77} and \textit{No. 39} anticipated legislative regulation of tenure, and the Decision of 1789 settled only that Congress may not reserve removal power to the Senate---not that Congress cannot impose for-cause limits. Sotomayor also pointed to founding-era semi-independent bodies (the Sinking Fund Commission, Mint Board, Revolutionary War Debt Commission, and First and Second Banks) and \textit{Marbury v. Madison}.

On \textit{stare decisis}, Sotomayor accused the majority of executing a ``profound bait and switch'' on Congress, ordinary Americans, and regulated businesses that relied on \textit{Humphrey's} to build dozens of bipartisan commissions---including FERC, the CPSC, the Chemical Safety Board, the NRC, and the MSPB (\href{https://www.law.cornell.edu/supremecourt/text/25-332}{Cornell LII}; \href{https://www.scotusblog.com/2026/06/court-allows-trump-to-fire-ftc-commissioner-and-overturns-major-restraint-on-presidential-power/}{SCOTUSblog}). Without removal protections, statutory bipartisanship, staggered terms, and multimember deliberation can be nullified overnight by firing opposition commissioners, concentrating unprecedented power in a single President.

\section{Examples: Future Applications}

\begin{enumerate}
    \item \textbf{Same-Side (At-Will Removal of an NLRB or FERC Commissioner):} Suppose a President fires a Democratic or Republican member of the National Labor Relations Board (29 U.S.C. \S~153) or the Federal Energy Regulatory Commission (42 U.S.C. \S~7171) purely over policy disagreements. Under \textit{Slaughter}, the removal is constitutional: both agencies promulgate substantive regulations, conduct binding enforcement adjudications, and pursue civil enforcement in federal court---falling squarely within the ``heartland of executive power'' where for-cause tenure protections are invalid.
    \item \textbf{Same-Side (Removal of a CPSC or NRC Commissioner):} Suppose a President summarily removes Commissioners of the Consumer Product Safety Commission or Nuclear Regulatory Commission because they favor stricter safety standards than the White House prefers. Because both commissions issue binding safety rules and enforce statutory mandates against private parties, \textit{Slaughter} renders their statutory for-cause removal restrictions unconstitutional, permitting at-will removal (\href{https://www.yalejreg.com/nc/implications-of-trump-v-slaughter-for-the-integrity-of-future-federal-commission-decisionmaking/}{Yale JREG}).
    \item \textbf{Opposite-Side (Purely Advisory or Congressionally Chartered Non-Executive Commission):} Suppose Congress creates a bipartisan Historical Preservation Advisory Commission tasked solely with researching federal historic sites and submitting nonbinding reports to Congress, with members removable only for cause. Under Part III--B of \textit{Slaughter} and \textit{Buckley v. Valeo}, the removal protection survives because the commission exercises only investigative and informative functions in aid of Congress and wields ``no part of the executive power.''
    \item \textbf{Opposite-Side (Federal Reserve Board of Governors):} Suppose a President attempts to fire a Federal Reserve Governor at will over interest-rate policy, citing \textit{Slaughter}. As Part III--B of \textit{Slaughter} explicitly signaled and \textit{Trump v. Cook} squarely held the same day, the Federal Reserve is an exception to \textit{Slaughter}'s general rule due to its unique historical lineage tracing to the First and Second Banks of the United States.
    \item \textbf{Fence-Sitter (Removal of a U.S. Tax Court Judge or Statutory Quorum / Party-Balance Challenges):} Suppose a President attempts to remove a Judge of the U.S. Tax Court without cause, notwithstanding 26 U.S.C. \S~7443(f), or fires all opposition-party commissioners at a commission with a fixed statutory quorum requirement and challenges the quorum and party-balance rules as unconstitutional burdens on Article II. Both scenarios fall squarely on the doctrinal frontier left open by \textit{Slaughter}: Part III--B expressly reserved the status of non-Article III courts (whose judges perform purely adjudicatory functions outside executive administration), and commentators such as former FERC Chair Joseph T. Kelliher have noted that \textit{Slaughter} leaves unresolved whether statutory quorum, qualification, and party-balance rules survive when they constrain presidential control (\href{https://www.yalejreg.com/nc/implications-of-trump-v-slaughter-for-the-integrity-of-future-federal-commission-decisionmaking/}{Yale JREG}).
\end{enumerate}

\section{Critique}

\textbf{Originalist and Textualist Views:} Proponents of the unitary executive hailed \textit{Slaughter} as a long-overdue restoration of Article II's original design, eliminating what President Roosevelt's own 1937 committee called a ``headless `fourth branch''' and replacing \textit{Humphrey's} fictional ``quasi'' labels with a clear, administrable rule of electoral accountability (\href{https://www.paulweiss.com/insights/client-memos/supreme-court-overrules-humphrey-s-executor-but-preserves-the-fed-s-independence}{Paul, Weiss}). Yet even sympathetic conservative scholars sounded serious alarms about the post-\textit{Slaughter} constitutional equilibrium. Writing in the \textit{Yale Journal on Regulation} (summarizing his \textit{Harvard Law Review} Supreme Court Comment), Professor Aaron L. Nielson identified the problem of ``severability stacking'': by repeatedly severing individual congressional checks---the legislative veto, for-cause removal, and oversight limits---while leaving intact massive statutory delegations of authority, the Court has created a ``statutory windfall'' of presidential power Congress never intended (\href{https://www.yalejreg.com/nc/the-american-presidency-post-slaughter/}{Yale JREG}). Similarly, Andy Smarick at \textit{Law \& Liberty} dubbed Justice Gorsuch's concurrence a ``con-jority opinion,'' observing that conservatives and liberals alike now share an urgent concern that handing cross-branch legislative and judicial powers to the President requires reviving nondelegation and Article III limits (\href{https://lawliberty.org/gorsuchs-con-jority-opinion-in-slaughter/}{Law \& Liberty}).

\textbf{Progressive, Historical, and Doctrinal Critiques:} Progressive scholars and historians---echoed by Justice Sotomayor's dissent---attacked the majority's historical narrative as selective and contradicted by recent archival scholarship showing that neither the English Crown, founding-era state constitutions, nor the Decision of 1789 barred Congress from limiting executive removals (\href{https://www.law.cornell.edu/supremecourt/text/25-332}{Cornell LII}). Doctrinally, critics faulted the Court for gutting \textit{stare decisis} despite nearly a century of congressional and private reliance, and for leaving a glaring inconsistency between \textit{Slaughter}'s absolutism and \textit{Cook}'s pragmatic Federal Reserve carveout. From a practical governance perspective, former FERC Chair Joseph T. Kelliher warned that \textit{Slaughter} risks triggering a ``Night of the Long Knives'' at every presidential transition, politicizing formerly merits-based regulatory commissions, deterring qualified experts from serving, and shattering long-term regulatory certainty for American businesses (\href{https://www.yalejreg.com/nc/implications-of-trump-v-slaughter-for-the-integrity-of-future-federal-commission-decisionmaking/}{Yale JREG}).

\textbf{Underlying Competing Values:} \textit{Slaughter} turns on a foundational clash between two visions of constitutional self-government: the majority's \emph{vertical accountability} model, in which every official executing federal law must answer to a single elected President so voters know whom to blame, versus the dissent's \emph{institutional pluralism and checks-and-balances} model, in which Congress may insulate bipartisan, expert commissions from short-term partisan whim to preserve regulatory stability, technical expertise, and impartial law enforcement.

\section{Key Quotations}

\begin{quote}
``When an agency `executes' a congressional mandate against private parties, it exercises executive power---no ifs, ands, or quasis about it.''
\end{quote}
\noindent --- Chief Justice Roberts, Opinion of the Court (\href{https://www.law.cornell.edu/supremecourt/text/25-332}{Cornell LII}; \href{https://www.supremecourt.gov/opinions/25pdf/25-322_new_4gcj.pdf}{Slip Op. at 22}).

\begin{quote}
``If anything more is left of \textit{Humphrey's}, we overrule it. \textit{Humphrey's} has for decades been a result in search of a rationale.''
\end{quote}
\noindent --- Chief Justice Roberts, Opinion of the Court (\href{https://www.law.cornell.edu/supremecourt/text/25-332}{Cornell LII}; \href{https://www.supremecourt.gov/opinions/25pdf/25-322_new_4gcj.pdf}{Slip Op. at 21}).

\begin{quote}
``Although it is up to the Senate to decide whether to confirm those with whom the President would prefer to work, neither Congress nor the courts may saddle him with those with whom he cannot work. Subordinates who exercise the President's power are subject to removal by him. Then, and only then, can they remain accountable to the President, and the President to the people.''
\end{quote}
\noindent --- Chief Justice Roberts, Opinion of the Court (\href{https://www.law.cornell.edu/supremecourt/text/25-332}{Cornell LII}; \href{https://www.supremecourt.gov/opinions/25pdf/25-322_new_4gcj.pdf}{Slip Op. at 36}).

\begin{quote}
``Open-ended delegations of legislative power have not gone away; now they will just be exercised by agency officials who answer to the President. The power to write new regulatory crimes still exists, but now the pen ultimately rests in the President's hand. The ability to judge disputes in-house remains, but now the house is white.''
\end{quote}
\noindent --- Justice Gorsuch, Concurring Opinion (\href{https://www.law.cornell.edu/supremecourt/text/25-332}{Cornell LII}; \href{https://supreme.justia.com/cases/federal/us/609/25-332/}{Justia}).

\begin{quote}
``In holding otherwise, the Court gives the President a power unknown even to the English Crown against which the Founders revolted, elevating him above his once-coequal branches by transforming a duty to take care that the laws be faithfully executed into a license to act in defiance of those very laws.''
\end{quote}
\noindent --- Justice Sotomayor, Dissenting Opinion (\href{https://www.law.cornell.edu/supremecourt/text/25-332}{Cornell LII}; \href{https://supreme.justia.com/cases/federal/us/609/25-332/}{Justia}).


\end{document}
